\begin{thebibliography}{99}
\bibitem[1]{ref1} Peter B. Andrews. An Introduction to Mathematical Logic and Type Theory: To Truth Through Proof, volume 27 of Applied Logic Series. Kluwer Academic Publishers, second edition, 2002. doi:10.1007/978-94-015-9934-4.
\bibitem[2]{ref2} Ali Assaf, Guillaume Burel, Raphaël Cauderlier, David Delahaye, Gilles Dowek, Catherine Dubois, Frédéric Gilbert, Pierre Halmagrand, Olivier Hermant, and Ronan Saillard. Dedukti: a Logical Framework based on the \ensuremath{\lambda}Π-calculus Modulo Theory, September 2016. Working Paper or Preprint. URL: https://hal.science/hal-01368161.
\bibitem[3]{ref3} Leo Bachmair and Harald Ganzinger. Rewrite-based equational theorem proving with selection and simplification. Journal of Logic and Computation, 4(3):217–247, 1994. doi:10.1093/ logcom/4.3.217.
\bibitem[4]{ref4} David Baelde, Kaustuv Chaudhuri, Andrew Gacek, Dale Miller, Gopalan Nadathur, Alwen Tiu, and Yuting Wang. Abella: A system for reasoning about relational specifications. J. of Formalized Reasoning, 7(2):1–89, 2014. doi:10.6092/issn.1972-5787/4650.
\bibitem[5]{ref5} David Baelde, Andrew Gacek, Dale Miller, Gopalan Nadathur, and Alwen Tiu. The Bedwyr system for model checking over syntactic expressions. In F. Pfenning, editor, 21th Conf. on Automated Deduction (CADE), number 4603 in LNAI, pages 391–397, New York, 2007. Springer. doi:10.1007/978-3-540-73595-3\_28.
\bibitem[6]{ref6} Gilles Barthe, Venanzio Capretta, and Olivier Pons. Setoids in type theory. Journal of Functional Programming, 13(2):261–293, 2003. doi:10.1017/S0956796802004501.
\bibitem[7]{ref7} Stanley Burris and H. P. Sankappanavar. A Course in Universal Algebra, volume 78 of Graduate Texts in Mathematics. Springer-Verlag, 1981. Millennium Edition available online. doi:10.1007/978-1-4613-8130-3.
\bibitem[8]{ref8} Alonzo Church. A formulation of the Simple Theory of Types. J. of Symbolic Logic, 5:56–68, 1940. doi:10.2307/2266170.
\bibitem[9]{ref9} Claudio Sacerdoti Coen. A semi-reflexive tactic for (sub-)equational reasoning. In Jean- Christophe Filliâtre, Christine Paulin-Mohring, and Benjamin Werner, editors, Types for Proofs and Programs (TYPES 2004), volume 3839 of Lecture Notes in Computer Science, pages 98–114. Springer, 2006. doi:10.1007/11617990\_7.
\bibitem[10]{ref10} Robert L. Constable et al. Implementing Mathematics with the Nuprl Proof Development System. Prentice-Hall, 1986.
\bibitem[11]{ref11} Gilles Dowek, Thérèse Hardin, and Claude Kirchner. Theorem proving modulo. J. of Automated Reasoning, 31(1):31–72, 2003.
\bibitem[12]{ref12} Cormac Flanagan, Amr Sabry, Bruce F. Duba, and Matthias Felleisen. The essence of compiling with continuations. ACM SIGPLAN Notices, 28(6):237–247, 1993. doi:10.1145/ 155090.155113.
\bibitem[13]{ref13} Peter J. Freyd. Properties invariant within equivalence types of categories. In Alex Heller and Myles Tierney, editors, Algebra, Topology, and Category Theory: A Collection of Papers in Honor of Samuel Eilenberg, pages 55–61. Academic Press, New York, 1976. doi:10.1016/ B978-0-12-339050-9.50010-0.
\bibitem[14]{ref14} Peter J. Freyd and Andre Scedrov. Categories, Allegories, volume 39 of North-Holland Mathematical Library. Elsevier Science Publishers B.V., Amsterdam, Netherlands, 1990. doi:10.1016/C2009-0-09144-1. 24:16   Treating Congruences as Equalities Within Proofs
\bibitem[15]{ref15} Andrew Gacek, Dale Miller, and Gopalan Nadathur. Nominal abstraction. Information and Computation, 209(1):48–73, 2011. doi:10.1016/j.ic.2010.09.004.
\bibitem[16]{ref16} Gerhard Gentzen. The consistency of elementary number theory. In M. E. Szabo, editor, The Collected Papers of Gerhard Gentzen, pages 132–213. North-Holland, Amsterdam, 1969. Translated from the 1936 German original. doi:10.1016/S0049-237X(08)70823-1.
\bibitem[17]{ref17} Ulysse Gérard and Dale Miller. Separating functional computation from relations. In Valentin Goranko and Mads Dam, editors, 26th EACSL Annual Conference on Computer Science Logic (CSL 2017), volume 82 of LIPIcs, pages 23:1–23:17, 2017. doi:10.4230/LIPIcs.CSL.2017.23.
\bibitem[18]{ref18} Jean-Yves Girard. A fixpoint theorem in linear logic. An email posting to linear@cs.stanford.edu archived at https://www.seas.upenn.edu/~sweirich/types/archive/1992/msg00030.html, February 1992.
\bibitem[19]{ref19} Joseph A. Goguen and José Meseguer. Equality, types, modules, and (why not?) generics for logic programming. The Journal of Logic Programming, 1(2):179–210, 1984. doi:10.1016/ 0743-1066(84)90011-X.
\bibitem[20]{ref20} Quentin Heath and Dale Miller. A proof theory for model checking. J. of Automated Reasoning, 63(4):857–885, 2019. doi:10.1007/s10817-018-9475-3.
\bibitem[21]{ref21} David Hilbert. The foundations of mathematics. In Jean van Heijenoort, editor, From Frege to Gödel: A Source Book in Mathematical Logic, 1879–1931, pages 464–479. Harvard University Press, Cambridge, Massachusetts, 1967. doi:10.4159/harvard.9780674324497.c22.
\bibitem[22]{ref22} Martin Hofmann and Thomas Streicher. The groupoid interpretation of type theory. In Twenty-Five Years of Constructive Type Theory, volume 36 of Oxford Logic Guides, pages 83–111. Oxford University Press, 1998. doi:10.1093/oso/9780198501275.003.0008.
\bibitem[23]{ref23} Oiva Ketonen. Investigations into the Predicate Calculus. College Publications, 2022. Ed. by S. Negri and J. von Plato.
\bibitem[24]{ref24} Laura Kovács and Andrei Voronkov. First-order theorem proving and Vampire. In Maria Paola Bonacina, editor, Automated Deduction - CADE-24 - 24th International Conference on Auto- mated Deduction, Lake Placid, NY, USA, June 9-14, 2013. Proceedings, volume 7898 of Lecture Notes in Computer Science, pages 1–35. Springer, 2013. doi:10.1007/978-3-642-38574-2\_1.
\bibitem[25]{ref25} Michael Makkai. First order logic with dependent sorts, with applications to category theory. Manuscript, McGill University, 1995. URL: https://www.math.mcgill.ca/makkai/folds/ foldsinpdf/FOLDS.pdf.
\bibitem[26]{ref26} Sonia Marin, Dale Miller, Elaine Pimentel, and Marco Volpe. From axioms to synthetic inference rules via focusing. Annals of Pure and Applied Logic, 173(5):1–32, 2022. doi: 10.1016/j.apal.2022.103091.
\bibitem[27]{ref27} Per Martin-Löf. Constructive mathematics and computer programming. In Sixth International Congress for Logic, Methodology, and Philosophy of Science, pages 153–175, Amsterdam, 1982. North-Holland.
\bibitem[28]{ref28} Jia Meng and Lawrence C. Paulson. Translating higher-order clauses to first-order clauses. Journal of Automated Reasoning, 40(1):35–60, 2008. doi:10.1007/s10817-007-9085-y.
\bibitem[29]{ref29} Dale Miller and Gopalan Nadathur. Programming with Higher-Order Logic. Cambridge University Press, June 2012. doi:10.1017/CBO9781139021326.
\bibitem[30]{ref30} Dale Miller and Alwen Tiu. A proof theory for generic judgments. ACM Trans. on Computa- tional Logic, 6(4):749–783, October 2005. doi:10.1145/1094622.1094628.
\bibitem[31]{ref31} John C. Mitchell and Gordon D. Plotkin. Abstract types have existential type. ACM Transactions on Programming Languages and Systems, 10(3):470–502, 1988. doi:10.1145/ 44501.45065.
\bibitem[32]{ref32} Sara Negri and Jan von Plato. Cut elimination in the presence of axioms. Bulletin of Symbolic Logic, 4(4):418–435, 1998. doi:10.2307/420956.
\bibitem[33]{ref33} Sara Negri and Jan von Plato. Structural Proof Theory. Cambridge University Press, 2001.
\bibitem[34]{ref34} Sara Negri and Jan von Plato. Proof Analysis: A Contribution to Hilbert’s Last Problem. Cambridge University Press, Cambridge, 2011. doi:10.1017/CBO9780511921629.
\bibitem[35]{ref35} Peter Schroeder-Heister. Rules of definitional reflection. In M. Vardi, editor, 8th Symp. on Logic in Computer Science, pages 222–232. IEEE Computer Society Press, IEEE, June 1993. doi:10.1109/LICS.1993.287585.
\bibitem[36]{ref36} Stephan Schulz, Simon Cruanes, and Petar Vukmirović. Faster, higher, stronger: E 2.3. 6th Vampire Workshop, 71:1–10, 2019. doi:10.29007/sc9t.
\bibitem[37]{ref37} Matthieu Sozeau. A new look at generalized rewriting in type theory. Journal of Formalized Reasoning, 2(1):41–62, 2009. doi:10.6092/issn.1972-5787/1574.
\end{thebibliography}
